% 问题三正文；由主论文与独立预览共同载入。
\section{问题三模型的建立与求解}
\label{sec:p3-model}

问题二通过接收约束筛选第二检测点，并用可能源位置下的交会区域直径评价定位效果。
本问将这一思路推广到多源连续观测：保留历次实测形成的可行域，在保证接收的候选点中，
兼顾缩域效果与移动耗时选择后续测点；同时加入未知频道搜索、清除判定与多源访问排序，
构成搜索、定位、清除交织执行的控制策略。

\subsection{任务状态与总时间目标}

沿用前文的目标圆域 $\Omega$、误差角域 $W(S,\theta)$ 和区域直径 $D(P)$。
记机器狗当前位置为 $X_t$、当前检测频道为 $b_t$，未知频道集合为 $\mathcal U_t$，
已发现及已清除频道集合分别为 $\mathcal F_t,\mathcal C_t$。对每个已发现源 $c$，维护其
位置可行多边形 $P_{c,k}$，其中 $k$ 为该频道已有的有效示向度数。初始时
$X_0=(0,0)$、$b_0=1$、$\mathcal U_0=\{1,\ldots,20\}$。

根据附录的计时规则，设总路程为 $L$，换频道、检测、清除失败及清除成功的次数依次为
$n_b,n_m,n_f,n_s$，则总虚拟时间为
\begin{equation}\label{eq:p3-time}
T=\frac{L}{5}+n_b+5n_m+3n_f+5n_s.
\end{equation}
移动与检测不能同时进行，清除指令不改变测向机频道。因此，式\eqref{eq:p3-time}
逐项累加实际动作成本，以全部源被清除为约束，尽量减小 $T$。未知源坐标无法预先读取，
故采用依赖当前观测的滚动决策，而不预设完整访问路线。程序现实运行时长另受模拟器
返回的剩余时限约束，不与虚拟时间混用。

\subsection{基于覆盖证书的未知频道搜索}

全向源的有效接收半径至少为 $1000$ m。记 $\mathcal S_t$ 为已完成当时全部未知频道扫描的
位置集合，定义其对场地的覆盖半径
\begin{equation}\label{eq:p3-cover}
R_{\rm cov}(\mathcal S_t)=\max_{G\in\Omega}\min_{S\in\mathcal S_t}\|G-S\|_2.
\end{equation}
若 $R_{\rm cov}\le1000$ m，则任何尚未发现的源都应至少在一个扫描点被检测到。
因此，仍未出现信号的频道可判为空；若已发现 $16$ 个不同频道的源，也可利用题设数量上限
结束未知频道搜索。单次无信号不能独立作为判空依据。

以圆心和半径 $r_0=1140$ m 的六个外圈点组成后备覆盖结构：
\begin{equation}\label{eq:p3-stations}
S_0=(0,0),\qquad
S_j=r_0\bigl(\cos(\phi+(j-1)\pi/3),\sin(\phi+(j-1)\pi/3)\bigr),\quad j=1,\ldots,6.
\end{equation}
其最大最近站距离约为 $992.689$ m，小于接收半径下限。计算一般扫描点集的覆盖半径时，
枚举场地内三站等距点、两站垂直平分线与场地圆周的交点，以及各站在圆周上的最远点，
并补入圆心和轴向边界点处理退化情形。对候选点计算最近站距离并取最大值，即得到
覆盖证书；其依据是最近站距离的最大值只能出现在相应 Voronoi 单元的内部顶点或圆周边界极值处。

定位与清除位置若完成未知频道扫描，也加入 $\mathcal S_t$。每轮枚举
$\phi=0^\circ,5^\circ,\ldots,55^\circ$，从相应六个后备点中选取补点集合 $\mathcal B_t$，
满足 $R_{\rm cov}(\mathcal S_t\cup\mathcal B_t)\le1000$ m。先选择补点数最少的方案，
再比较其与已知源区域联合规划的路线长度。后备点最多六个，每个角度至多检查 $2^6$ 个子集。
这一选择用于减少后续搜索负担，并不意味着获得全局时间最优解。

\subsection{保留历史观测的源位置可行域}

问题二首次观测的源可行域综合了角域、场地及距离信息，本问在此基础上保留全部历史观测。
记 $\overline\Omega$ 和 $\overline B(S,1500)$ 分别为场地圆盘与接收上界圆盘的外接正180边形，
首次发现前取 $P_{c,0}=\overline\Omega$。在 $S$ 获得示向度 $\theta$ 后，令
$u_\theta=(\cos\theta,\sin\theta)$，按以下方式更新：
\begin{equation}\label{eq:p3-update}
P_{c,k+1}=P_{c,k}\cap W(S,\theta)\cap\overline B(S,1500)
\cap\{G:u_\theta\cdot(G-S)\ge5\cos\alpha\},\quad \alpha=1^\circ.
\end{equation}
最后一项由正常示向度对应距离大于5 m得到，是距离下界的保守线性替代。
外接多边形包含真实圆盘，各约束可通过半平面裁剪求交；在观测误差满足所用角界的条件下，
更新后区域仍包含真实源。此处使用含距离信息的可行域，而非仅由角域相交得到的区域。

全向源的固定接收半径还使无信号记录具有定位价值。设同一未清除源在 $S$ 点有信号、
在 $N$ 点无信号，则 $\|G-S\|\le R_c<\|G-N\|$。平方展开可得
\begin{equation}\label{eq:p3-nosignal}
2(N-S)\cdot G\le\|N\|_2^2-\|S\|_2^2.
\end{equation}
用闭半平面保守保留严格不等式的边界，与 $P_{c,k}$ 求交即可进一步缩域。发现源时调用
此前保存的无信号记录；以后获得无信号时则与已有有效测向点配对。若返回近场信息
\texttt{near}，已知距离不超过5 m，直接在当前位置清除。

区域直径按问题一的顶点最远距离结论计算，程序采用顶点对枚举，顶点数为 $h$ 时
复杂度为 $O(h^2)$。另需区分角误差与读数舍入：当前实现的角界为
$1^\circ$，接口示向度保留两位小数；若按先加误差再舍入处理，严格外包应计入至多
$0.005^\circ$ 的舍入量。正式验证时需核验该边界条件，清除余量不能替代角界校核。

\subsection{由问题二选址模型到时间型主动测向}

\textbf{接收约束与效果预测。}
问题二采用“接收可行性筛选、可能源位置采样、近优点保留”的选址框架，本问沿用这一框架，
并针对当前多边形检查全部顶点。省略下标，记 $P=P_{c,k}$、顶点集为 $V(P)$，主动测点集合为
\begin{equation}\label{eq:p3-receive}
\mathcal A_{c,k}=\{q:\max_{v\in V(P)}\|q-v\|_2\le1000\}.
\end{equation}
由凸性，该条件覆盖多边形内全部可能位置，是与源实际接收半径无关的充分条件。
测点不必受限于目标圆域，附录允许机器狗在场地外检测；程序对部分候选另作场地筛选。

从 $P$ 取顶点、部分边中点及顶点均值组成样本集 $\mathcal G$。设 $\beta(q,G)$ 为 $q$ 指向
$G$ 的方位角，以 $E=\{-1^\circ,0,1^\circ\}$ 模拟读数误差，定义
\begin{equation}\label{eq:p3-predict}
\widehat D_c(q)=\max_{G\in\mathcal G,\ e\in E}
D\bigl(P\cap W(q,\beta(q,G)+e)\bigr).
\end{equation}
对 $\|q-G\|\le5$ m 的近场样本，将其后续定位需求记为零。与问题二的名义交会评价相比，
式\eqref{eq:p3-predict}固定已有实测形成的 $P$，仅模拟下一次观测，因而能逐次使用新的定位信息。
有限样本及三个误差值仅用于选点评价，不代表连续源位置和误差下的严格最坏直径。

\textbf{时间评分与粗细网格求解。}
设当前最小包围圆圆心为 $O_c$，程序对测点采用
\begin{equation}\label{eq:p3-score}
\widehat T_c(q)=\frac{\|X_t-q\|_2}{5}+5
+w_k\frac{\|q-O_c\|_2}{5}+0.85\frac{\widehat D_c(q)}{5},
\quad
w_k=\begin{cases}1.05,&k\le2,\\0.38,&k\ge3.\end{cases}
\end{equation}
前两项为到达与检测成本，后两项估计测后接近源及继续定位的代价；跨频道比较动作时再加换台耗时。
早期较大的圆心距离权重鼓励向源区域推进。该评分是时间导向的启发式量，不等于最终完成时间。

首先生成125 m粗网格，筛除不满足式\eqref{eq:p3-receive}的点，并优先保持与历史测点至少
80 m的基线，无候选时放宽至1 m。利用检测距离与交会角构造几何指标粗筛，同时补入靠近
当前位置、圆心及接近路径的候选，再按式\eqref{eq:p3-predict}精评估。记本轮最好预测值为
$\widehat D_{\min}$，保留 $\widehat D_c(q)\le\eta_k\widehat D_{\min}$ 的点，
其中 $k\le2$ 时 $\eta_k=1.50$，否则取1.18。在该近优集合中最小化式\eqref{eq:p3-score}，
再以所选粗点为中心，用30 m网格局部精化。由此将问题二的定位质量筛选转化为兼顾精度与路程的连续选点策略。

\subsection{单点清除与局部扫清决策}

由问题一的圆覆盖结论，区域直径的一半不能直接视为覆盖半径。对 $P_{c,k}$ 求最小包围圆：
\begin{equation}\label{eq:p3-mec}
r_c=\min_O\max_{v\in V(P_{c,k})}\|O-v\|_2,\qquad O_c\in\arg\min_O\max_v\|O-v\|_2.
\end{equation}
若 $r_c\le19$ m，则在 $O_c$ 处到任一可能源位置的距离均小于20 m，可执行单点清除，
并以接口成功反馈更新清除状态。19 m阈值为距离计算保留余量。

若区域已较小但不满足上述条件，可尝试局部扫清。对多边形取若干旋转方向的外接矩形，
划分网格，使每格宽、高 $a,b$ 满足 $\sqrt{a^2+b^2}/2\le18.5$ m。各格中心的清除圆
覆盖整个矩形，从而覆盖 $P_{c,k}$。只接受至多6个清除点的方案，按短路线逐点尝试，成功即停止。
若有 $m$ 个点、含到达首点的总路程为 $L_g$，则完整扫清的最坏耗时为
\begin{equation}\label{eq:p3-sweep}
T_g=L_g/5+3(m-1)+5.
\end{equation}
程序在 $T_g\le\widehat T_c(q)+\mathbf1_{b_t\ne c}+45$ s 时选择扫清，否则继续测向。
45 s为对后续定位成本的经验补偿。扫清保证指整个覆盖序列至少一次成功，并非每次尝试都成功。

\subsection{多任务位置复用与滚动路线求解}

将未清除源的 $O_c$ 与后备覆盖点 $\mathcal B_t$ 组成节点集合，从 $X_t$ 出发用最近邻法
生成初始顺序，再以至多8轮2-opt反转路段、缩短路程。任务无需回到原点，采用开放路线。
源圆心仅用于区域排序，实际处理该源时仍按前两节选择测点或清除动作。每轮只处理路线
首个节点，随后更新信息并重新规划。

在覆盖点扫描未知频道后，对有缩域价值的已知源补测；若不能保证接收，则评价无信号
半平面可能带来的收缩。定位或清除位置距已有扫描点至少1200 m时，扫描未知频道并补入覆盖集合；
其余已知源在满足接收条件且预测缩域充分时，可原地补测，每个普通停靠点至多补测5个源。
额外扫描仍计入式\eqref{eq:p3-time}，位置复用节省的是专程移动。

图\ref{fig:p3-flow}给出整体流程。路线不可用时，以当前动作代价、下一任务距离和同点补测收益
进行备用排序。完整任务的结束条件为未知频道已判定完毕且所有发现源均确认清除；
计算与通信均需在现实截止时间内完成，动作请求串行发送。

\begin{figure}[H]
\centering
% 紧凑矢量流程图；只依赖 tikz 与 arrows.meta。
\begin{tikzpicture}[x=1cm,y=1cm,>=Latex,font=\small,
  p3box/.style={draw,rounded corners=2pt,align=center,minimum height=0.63cm,inner sep=4pt},
  p3arrow/.style={->,line width=0.55pt}]
\node[p3box,text width=4.8cm] (start) at (0,0) {进入模拟器，原点扫描未知频道};
\node[p3box,text width=4.8cm] (check) at (0,-0.95) {更新可行域、覆盖与清除状态};
\node[p3box,text width=2.8cm] (end) at (4.7,-0.95) {退出并记录结果};
\node[p3box,text width=4.8cm] (plan) at (0,-1.9) {补齐覆盖候选，联合规划开放路线};
\node[p3box,text width=3.8cm] (scan) at (-1.7,-2.95) {覆盖节点：扫描未知频道};
\node[p3box,text width=4.3cm] (act) at (3.3,-2.95) {源节点：测向／单点清除／扫清};
\node[p3box,text width=4.8cm] (reuse) at (0,-4.0) {按条件复用位置，扫描或补测};
\draw[p3arrow] (start)--(check);
\draw[p3arrow] (check)--node[above,font=\footnotesize]{全部完成}(end);
\draw[p3arrow] (check)--node[right,font=\footnotesize]{未完成}(plan);
\draw[p3arrow] (plan.south)--++(0,-0.15)-|(scan.north);
\draw[p3arrow] (plan.south)--++(0,-0.15)-|(act.north);
\draw[p3arrow] (scan.south)--++(0,-0.17)-|(reuse.north);
\draw[p3arrow] (act.south)--++(0,-0.17)-|(reuse.north);
\draw[p3arrow] (reuse.west)--(-4.3,-4.0)--(-4.3,-0.95)--(check.west);
\end{tikzpicture}

\caption{搜索、定位与清除的滚动决策流程。每轮处理一个节点，动作串行执行。}
\label{fig:p3-flow}
\end{figure}

\newpage
\subsection{测试安排与结果记录}

通过不限次数的演练检查覆盖搜索、可行域更新和清除反馈，再按题目要求完成三次正式测试。
清除比例与平均定位清除时间分别为
\begin{equation}\label{eq:p3-metrics}
\rho_{\rm clear}=\frac{N_{\rm cleared}}{N_{\rm true}},\qquad
\overline T=\frac{T}{N_{\rm cleared}}.
\end{equation}
$N_{\rm true}$ 应取模拟器提供的实际源数，不能用程序发现数代替；无成功清除时不计算平均值。
按题目表1记录正式结果，如表\ref{tab:p3-official}所示。

\begin{table}[H]
\centering
\caption{问题3正式测试结果（按题目表1格式，待正式测试后填写）}
\label{tab:p3-official}
\small
\begin{tabularx}{\textwidth}{@{}cYYYY@{}}
\toprule
 & 测试案例编码 & 清除干扰源个数 & 平均定位清除时间/s & 程序运行时间/s\\
\midrule
测试1 & 待填 & 待填 & 待填 & 待填\\
测试2 & 待填 & 待填 & 待填 & 待填\\
测试3 & 待填 & 待填 & 待填 & 待填\\
\bottomrule
\end{tabularx}
\end{table}

正式测试数据尚待补充，表中仅保留填写位置，暂不作数值性能结论。后续结合完整日志核对
式\eqref{eq:p3-time}各项成本，分析路程与测量次数对耗时的贡献，并保存原文件名的三次正式
测试日志作为支撑材料。搜索覆盖和清除覆盖的几何依据，与选点及路线的启发式效果应分开检验；
近优门槛本身不证明有限步收敛，程序循环上限也仅是异常保护，仍需检查困难案例及现实耗时。
